\section{フレーム依存性と適用範囲}
\subsection{初期シアーなしフレームの必要性}
$u$ に依存しないシアーのずれ
\[
C_{AB}\mapsto C_{AB}+S_{AB}
\]
ではニュースは変わらないが、ハード汎関数は
\[
\bar\Phi_{\rm H}\mapsto\bar\Phi_{\rm H}
-\frac18\Delta C_{AB}\widetilde S^{AB}
\]
だけ変化する。一方、式\eqref{eq:area}の弦で閉じた面積は $(p,q)\mapsto(p+p_0,q+q_0)$ に不変である。したがって、Jacobi 面積とハード項単独の同一視は一般の一定シアー基準では成立しない。初期シアーなしフレームは単なる記法上の便宜ではなく、ハード成分だけの対応関係を成立させる物理的規約である。

\subsection{固定ヌル窓が補助構造であること}
$Q_\gamma$ は波形だけから決まる固有量ではなく、固定したヌル窓 $\gamma$ を含む補助的なコストである。例えば曲率の台と放射履歴を変えずに、両端へ曲率ゼロ・ニュースゼロの定常区間だけを付加し、同じ $du/d\lambda=1$ 規格化のまま窓長を $L$ から $L'>L$ へ延長する。このとき $\int_\gamma B(k^4)d\lambda$ とハード・メモリーは変わらない一方、
\[
Q_{\gamma'}=\left(\frac{L'}{L}\right)^3Q_\gamma
\]
となる。従って同じ物理波形について窓を延長すれば $Q_\gamma$ を任意に大きくできる。本稿の最小作用則は、固定窓ごとの無次元化された変分問題に対するものであり、異なる窓の $Q_\gamma$ を同一の detector exposure として比較する主張ではない。なお時間軸と波形全体を同時にアフィン再パラメータ化する場合は、第2節で示した通り $Q_\gamma$ は不変であり、これは空の定常尾を追加して物理的な窓そのものを変える操作とは異なる。

\subsection{定常条件が最良係数を変える理由}
無拘束最適化では $K_0$ の最大特異値部分空間が回復族を与えるが、定常な放射では $\int\alpha=\int\beta=0$ が必要である。この条件は各偏光に余次元 1 の閉部分空間を課すため、稠密近似によって無視できない。実際、全空間の最大特異値部分空間に $\int\beta=0$ と $\int K_0\beta=0$ を同時に課す線形系は数値的に非退化であり、旧最適化波形は定常クラスに入らない。正しい接線方向作用素は $K_{\rm stat}=P_0K_0P_0$ であり、$\|K_{\rm stat}\|<\|K_0\|$ なので、必要な鋭い曲率曝露係数は増加する。

本稿は最終シアーを零には固定しない。もしさらに $p(1)=q(1)=0$ を課すなら、平均ゼロ条件に加えて
\[
\int_0^1(1-t)\alpha(t)\,dt=0,\qquad
\int_0^1(1-t)\beta(t)\,dt=0
\]
という一次モーメント制約が加わる。その場合の許容部分空間と圧縮作用素は $K_{\rm stat}$ とは異なり、最良係数も一般には変わりうるため、本稿の $71.13876223\ldots$ をその追加制約付き問題へ移してはならない。

\subsection{主張しないもの}
本稿は次を主張しない。
\begin{itemize}
\item ジャイロスコピック・メモリー全体への鋭い下界；
\item スピン・メモリーへの鋭い下界；
\item 一般の BMS フレームで不変なハード成分のみの定理；
\item 有限振幅での大域的最小作用則；
\item 本稿の回復 TT/Bondi データが任意に完全な漸近平坦真空 Einstein 時空へ大域実現できるという主張；
\item 有限 $r$ での全次数の厳密恒等式；
\item Bondi エネルギー、Isaacson エネルギー、放射エネルギー、検出器エネルギーの下界；
\item ジャイロスコピック・メモリー、Bel--Robinson テンソル、Jacobi メモリーテンソル自体の発見。
\end{itemize}

従って主定理\ref{thm:main-physical}の「物理クラス」は、完全 Einstein 解の大域解空間ではなく、標準的な初期シアーなし規約のもとで整合する弱振幅・放射域先頭次数の TT/Bondi 放射データクラスを意味する。比較している量も full finite-amplitude の $Q_\gamma$ と $\Delta\Psi_{\rm H}$ ではなく、その二次摂動係数 $Q_\gamma^{(2)}$ と $\Delta\Psi_{\rm H}^{(2)}$ である。完全な非線形時空への実現可能性と有限振幅での値関数は別問題として残す。

\section{数値検証と再現性}
解析証明と数値検証を区別する。式\eqref{eq:sharp-jacobi},\eqref{eq:stationary-sharp},\eqref{eq:bridge},\eqref{eq:main-physical-value}の論理自体は解析的であり、以下の数値は規格化、符号、終端量の抽出、制限スペクトル値、および反例志向の検査例を確認するためのものである。

\subsection{\texorpdfstring{無拘束 $K_0$ と厳密 Jacobi 回復}{無拘束 K0 と厳密 Jacobi 回復}}
Gauss--Legendre Nystr\"om 離散化は
\begin{center}
\begin{tabular}{rr}
\toprule
節点数 & $\|K_0\|$\\
\midrule
50&0.0910384416041\\
100&0.0910384384527\\
200&0.0910384382516\\
300&0.0910384382407\\
500&0.0910384382384\\
\bottomrule
\end{tabular}
\end{center}
を与える。同じ離散最大特異値対を厳密 Jacobi ODE に入れると
\begin{center}
\begin{tabular}{rr}
\toprule
$\rho$ & $Q_\gamma/|\omega_\gamma|$\\
\midrule
0.03&10.98435328\\
0.01&10.98436922\\
0.003&10.98437103\\
0.001&10.98437118\\
\bottomrule
\end{tabular}
\end{center}
となり、$\|K_0\|^{-1}=10.98437121\ldots$ に収束する。ここで $\rho$ は回復半直線 $A_\rho=\rho A_*$ の振幅であり、物理摂動パラメータ $\eps$ とは別である。反対称 Jacobi チャネルでは奇数次数が消えるため $\omega_\gamma[A_\rho]=2s_0\rho^2+O(\rho^4)$ となり、比の誤差が $O(\rho^2)$ で収束することとも整合する。

\subsection{定常条件付きスペクトルと厳密回復}
平均ゼロ射影を離散化して $P_0MP_0$ の最大特異値を計算すると、
\begin{center}
\begin{tabular}{rrr}
\toprule
節点数&$\|K_{\rm stat}\|$&$2/\|K_{\rm stat}\|$\\
\midrule
50&0.0281140528760&71.13880054294\\
100&0.0281140670547&71.13876466581\\
200&0.0281140679578&71.13876238062\\
300&0.0281140680065&71.13876225721\\
500&0.0281140680170&71.13876223067\\
\bottomrule
\end{tabular}
\end{center}
を得る。さらに連続な平均ゼロ最大特異値対を厳密 Jacobi ODE へ入れると
\begin{center}
\begin{tabular}{rr}
\toprule
$\rho$ & $Q_\gamma/|\omega_\gamma|$\\
\midrule
0.03&35.56942171300\\
0.01&35.56938562391\\
0.003&35.56938151880\\
0.001&35.56938115793\\
\bottomrule
\end{tabular}
\end{center}
となり、$\|K_{\rm stat}\|^{-1}=35.5693811\ldots$ へ収束する。ここでも比の誤差は上記の偶数次数構造により $O(\rho^2)$ である。再構成したひずみは数値精度で $p'(1)=q'(1)=0$ を満たす。連続最大特異値対を $\|\alpha_*\|=\|\beta_*\|=1$ と規格化すると、数値 SVD の一つの向き付けでは $p_f/\rho\simeq-0.47863$, $q_f/\rho\simeq0.30730$ である。最大特異ベクトル対には全体符号の任意性があるため、両者を同時に反転した $(0.47863,-0.30730)$ も同じ最適化器を表す。偏光基底の回転とこの全体符号に依存しない最終シアーの大きさは $\sqrt{p_f^2+q_f^2}/\rho\simeq0.56878$ である。本定理は最終シアーを零に固定せず、一定の displacement-memory 成分を許している。これらは制限定理そのものの証明ではなく、解析的回復構成の実装検査である。

\subsection{対応関係の波形検査}
$W=\int(p\dot q-q\dot p)du$ とすると、初期ゼロ基準の場合 $\omega^{(2)}=W/2$, $\Delta\Psi_{\rm H}^{(2)}=-W/4$ である。代表的な検査例は
\begin{center}
\begin{tabular}{lrrrr}
\toprule
波形族&$W$&$\omega^{(2)}$&$\Delta\Psi_{\rm H}^{(2)}$&比\\
\midrule
滑らかな円偏光&2.35619447&1.17809724&-0.58904862&-2.000000\\
逆向きヘリシティ&-2.35619447&-1.17809724&0.58904862&-2.000000\\
楕円偏光&1.22522112&0.61261056&-0.30630528&-2.000000\\
固定軸&$\simeq0$&$\simeq0$&$\simeq0$&ともにゼロ\\
終端メモリー&-0.62590250&-0.31295125&0.15647563&-2.000000\\
ヘリシティ相殺&$\simeq0$&$\simeq0$&$\simeq0$&ともにゼロ\\
\bottomrule
\end{tabular}
\end{center}
である。基準値に対する反例志向の検査では、終端メモリー波形に $(p_0,q_0)=(0.35,-0.22)$ を加える前後で、弦で閉じた面積は $-0.3129512514\to-0.3129512515$ と不変だが、ハード・ジャイロ汎関数は $0.1564756257\to0.1124756258$ と変化する。これは一般の基準値に対するハード成分のみの恒等式を反証し、初期シアーなしフレーム制限を支持する。

\subsection{旧最適化波形の不適合を検出する監査}
全空間の $K_0$ 最大特異値部分空間に正規直交基底を取り、$\int\beta=0$ と $\int K_0\beta=0$ の二つの制約行を $\|K_0\|$ で規格化すると、その $2\times2$ 行列式の絶対値は基底の直交変換では不変で、$n=500$ で約 $0.90925$ となる。この量は同時に、定数関数の最大特異部分空間への直交射影のノルム二乗に等しい。従って定数モードの約 $90.9\%$ が旧最大特異平面に含まれ、二つの定常制約がその平面上で独立であることを示す。これは旧無拘束最大特異値部分空間が定常条件を自動的には満たさないことを示す。この監査は「なぜ旧係数が鋭くなかったか」を診断するためのものであり、新係数の鋭さ自体は定理\ref{thm:stationary-sharp}の解析証明から従う。
